\documentclass[10.5pt]{article}

\usepackage[a4paper,margin=0.7in]{geometry}
\usepackage{titlesec}
\usepackage{enumitem}
\usepackage{hyperref}

\setlength{\parindent}{0pt}

\hypersetup{
    colorlinks=true,
    linkcolor=black,
    urlcolor=blue
}

\setlist[itemize]{noitemsep, topsep=0pt, leftmargin=*}
\titleformat{\section}{\small\bfseries}{}{0em}{}[\titlerule]

\begin{document}

% -------------------- Header --------------------
\begin{center}
    {\Large \textbf{Anushka Dhiman}} \\[4pt]
    \href{mailto:anushkadhiman101@gmail.com}{anushkadhiman101@gmail.com} \,|\, 
    9027996957 \,|\, 
    \href{https://www.linkedin.com/in/anushka-dhiman-7b306027b/}{LinkedIn} \,|\, 
    \href{https://github.com/anushkaaa26}{GitHub}
\end{center}

% -------------------- Education --------------------
\section*{Education}
\textbf{VIT Bhopal University}, Bhopal \hfill Sep 2023 -- Jun 2027 \\
B.Tech in Computer Science (AI \& ML) \\
\textbf{Delhi Public School}, Meerut \\
Class X and XII

% -------------------- Experience --------------------
\section*{Experience}

\textbf{Infosys Springboard} — AI/ML Intern (Remote) \hfill Nov 2025 -- Dec 2025
\begin{itemize}
    \item Built \textbf{AI-ImpactSense}, an ML-based earthquake impact prediction system using supervised learning models.
    \item Performed data preprocessing, feature engineering, and dataset structuring to improve model reliability.
    \item Implemented end-to-end ML pipeline including training, evaluation, and hyperparameter tuning.
    \item Deployed model using \textbf{Streamlit} with real-time predictions and risk visualization.
\end{itemize}

\textbf{Virsis} — Project Intern (Remote) \hfill Jan 2025 -- Sep 2025
\begin{itemize}
    \item Improved organic traffic through SEO strategies and keyword optimization.
    \item Analyzed web traffic data and generated performance reports with actionable insights.
\end{itemize}

% -------------------- Skills --------------------
\section*{Skills}
\textbf{Programming:} Python \\
\textbf{Computer Vision:} Object Detection (YOLOv8), Image Classification, OCR (Tesseract), Image Processing \\
\textbf{Machine Learning:} Supervised Learning, Deep Learning, Model Optimization, Feature Engineering \\
\textbf{Frameworks:} PyTorch, TensorFlow, Keras, OpenCV, Scikit-learn \\
\textbf{Data \& Pipelines:} Data Preprocessing, Data Augmentation, Dataset Annotation, Pipeline Development \\
\textbf{Systems \& Deployment:} Real-time Systems, Multi-module Architecture, Streamlit, Flask APIs \\
\textbf{Concepts:} Computer Vision Pipelines, Model Evaluation, Hyperparameter Tuning, Event-driven Systems

% -------------------- Projects --------------------
\section*{Projects}

\textbf{Netra-AI | Real-Time Assistive Navigation System} \hfill
\href{https://github.com/anushkaaa26/NetraAI}{GitHub} \\
\begin{itemize}
    \item Designed and developed an end-to-end \textbf{computer vision system} integrating multiple real-time modules.
    \item Implemented \textbf{YOLOv8-based object detection} achieving $\sim$90\% accuracy at 50--90ms/frame.
    \item Built custom-trained \textbf{YOLO model} for Indian currency recognition with $\sim$87\% accuracy.
    \item Integrated OCR (Tesseract) and image processing pipelines for real-time text extraction.
    \item Engineered a multi-module pipeline including obstacle detection, fall detection, OCR, and navigation.
    \item Developed BFS-based navigation system and voice-driven interaction for hands-free operation.
\end{itemize}
\textbf{Tech Stack:} Python, PyTorch, OpenCV, YOLOv8, Tesseract OCR

\vspace{2pt}

\textbf{AI-ImpactSense | Earthquake Impact Prediction System} \hfill
\href{https://ai-impactsense-jghxyvcnxdftwkqjgr3imr.streamlit.app/}{Website} \,|\, 
\href{https://github.com/anushkaaa26}{GitHub} \\
\begin{itemize}
    \item Developed ML models to predict earthquake impact severity with structured datasets.
    \item Performed data cleaning, preprocessing, and feature engineering for improved performance.
    \item Built a complete ML pipeline with evaluation metrics and visualization.
\end{itemize}
\textbf{Tech Stack:} Python, Pandas, Scikit-learn, Matplotlib, Streamlit

\vspace{2pt}

\textbf{Dog Breed Classifier | Deep Learning} \hfill
\href{https://github.com/anushkaaa26/dog_breed_identification}{GitHub} \\
\begin{itemize}
    \item Built a multi-class image classification model using CNNs with TensorFlow/Keras.
    \item Applied image preprocessing and data augmentation to improve generalization.
    \item Deployed model for real-time inference via web interface.
\end{itemize}
\textbf{Tech Stack:} Python, TensorFlow, Keras, OpenCV

% -------------------- Certifications --------------------
\section*{Certifications}
\begin{itemize}
    \item Applied Machine Learning in Python — University of Michigan (Nov 2025)
    \item Job Simulation — Forage Academy (2024)
\end{itemize}

\end{document}
